\subsubsection{Logical vs Physical Qubits}\label{sec:logical-vs-physical-qubits}

\textcolor{Red2}{\faIcon{exclamation-triangle} \textbf{Problem.}} In an \hl{ideal mathematical model of quantum computation, qubits are assumed to evolve perfectly under unitary transformations.} Real quantum hardware, however, is inherently \textbf{noisy}. Physical qubits interact with the surrounding environment and are affected by decoherence, thermal fluctuations, imperfect control pulses, and measurement errors. As a consequence, \textbf{quantum information} stored in a single physical qubit is \textbf{extremely fragile}.

\highspace
\textcolor{Green3}{\faIcon{check-circle} \textbf{Solution.}} To overcome this limitation, quantum computers do not directly perform large computations on individual physical qubits. Instead, they use \textbf{quantum error correction}, where the information of a single qubit is encoded across many physical qubits, forming a \textbf{logical qubit}.
\begin{itemize}
    \item \definition{Physical Qubit}. A \textbf{physical qubit} is the actual \textbf{hardware-level quantum system used to store information}. For example superconducting qubits, trapped ions, neutral atoms or photonic qubits\footnote{%
        It's normal don't know the details of these different physical implementations at this stage. The key point is that they are all examples of physical qubits, which are the building blocks of quantum computers.
    }. Physical qubits are \textbf{directly manipulated by the hardware} and \textbf{implement the elementary quantum operations supported by the quantum processor}. However, they are noisy and have limited coherence times.


    \item \definition{Logical Qubit}. A \textbf{logical qubit} is an \textbf{encoded qubit constructed from many physical qubits} in order \hl{to protect quantum information against errors.}

    Instead of storing the quantum state in a single fragile qubit, the \hl{information is distributed redundantly across multiple physical qubits:}
    \begin{equation*}
        \text{Many physical qubits} \xrightarrow{\text{Encoding}} \text{One logical qubit}
    \end{equation*}
    This encoding allows the system to detect errors, correct certain errors and preserve the quantum state for a much longer time.

    \textcolor{Green3}{\faIcon{question-circle} \textbf{Why do we need logical qubits?}} 
    Without error correction, errors accumulate extremely rapidly during a quantum computation. Even a small error probability per gate becomes catastrophic when millions or billions of gates are executed.

    Logical qubits make large-scale quantum computation possible by continuously protecting the encoded information while the computation is running.
\end{itemize}
\begin{flushleft}
    \textcolor{Red2}{\faIcon{exclamation-triangle} \textbf{The Main Challenge}}
\end{flushleft}
\hl{Encoding a logical qubit is only the first step.} We must also be able to apply quantum gates, manipulate information, execute algorithms. Directly on logical qubits without destroying the error-correcting structure. This leads to the concept of \textbf{fault-tolerant quantum gates}, which are \hl{quantum gates specifically designed to prevent local physical errors from spreading uncontrollably through the encoded system.}